\chapter{Testing and Results}

This chapter reports the verification evidence for NEXIS as it actually
occurred during the final engineering cycle. Rather than presenting a
constructed matrix of numbered test cases, the chapter states three kinds of
real evidence: (i) static verification --- the build, vet, lint, type-check,
and automated test commands that were run against every service, and their
results; (ii) live end-to-end integration runs performed against the real,
running system --- two complete nine-agent recovery pipeline executions and
one complete deploy-engine execution, all driven by a real local large
language model rather than stubs or mocks; and (iii) the defects that this
testing process itself uncovered, together with their root causes and fixes.
The chapter closes with an honest summary of what was and was not tested.

\section{Static Verification}

Every component of the platform was verified with its native toolchain. For
the four independent Go modules, the full sequence \code{go build ./...},
\code{go vet ./...} (Go's built-in static analyser), and \code{go test ./...}
was executed and required to complete with zero errors, zero vet findings,
and zero test failures. The two Python components were verified with
\code{pytest}; the causal-inference service's suite comprises 38 tests, and
the validator's Hypothesis-based sidecar~\cite{hypothesis} has its own
passing suite. The Next.js frontend was verified with the project's three
gate commands: TypeScript type checking, ESLint, and a full production
build. Table~\ref{tab:static-verification} summarises the commands and
their outcomes.

\begin{table}[htbp]
\caption{Static verification commands and results across all NEXIS components.}
\label{tab:static-verification}
\centering
\small
\begin{tabular}{@{}p{3.5cm}p{6.3cm}p{4.6cm}@{}}
\toprule
\textbf{Component} & \textbf{Command} & \textbf{Result} \\
\midrule
\code{services/control-plane} (Go) & \code{go build ./... \&\& go vet ./... \&\& go test ./...} & Clean --- build, vet, and all tests pass \\
\addlinespace
\code{services/validator} (Go) & \code{go build ./... \&\& go vet ./... \&\& go test ./...} & Clean --- build, vet, and all tests pass \\
\addlinespace
\code{services/gitops} (Go) & \code{go build ./... \&\& go vet ./... \&\& go test ./...} & Clean --- build, vet, and all tests pass \\
\addlinespace
\code{services/deploy-engine} (Go) & \code{go build ./... \&\& go vet ./... \&\& go test ./...} & Clean --- build, vet, and all tests pass \\
\addlinespace
\code{services/causal-inference} (Python) & \code{pytest} & 38 tests, all passing \\
\addlinespace
Validator Hypothesis sidecar (Python) & \code{pytest} & All tests passing \\
\addlinespace
\code{apps/web} (Next.js) & \code{pnpm typecheck \&\& pnpm lint \&\& pnpm build} & Clean --- type check, lint, and production build all pass \\
\bottomrule
\end{tabular}
\end{table}

Two aspects of this table deserve emphasis. First, the result column is not
aspirational: every command listed was actually executed against the final
codebase and completed cleanly, and the Go test suites include the JWT
clock-injection tests discussed in Section~\ref{sec:defects}, which were
silently failing earlier in the cycle and now genuinely pass. Second, static
verification covers every deployable unit of the platform --- there is no
service that ships without passing its build, its static analysis, and its
automated tests.

\section{Live End-to-End Recovery Runs}

Static verification establishes that each service is internally sound; it
says nothing about whether the nine agents, the Temporal
workflow~\cite{temporal}, the validator sandbox, the approval gate, and the
live console actually cooperate. For that, two complete recovery pipeline
runs were executed against the real running system --- real Postgres with
row-level security active, real Temporal workflows, real Neo4j dependency
graph, and a real local LLM (Ollama~\cite{ollama} serving
\code{qwen2.5-coder:14b} and \code{gpt-oss:20b}) generating the actual
plans, diffs, and tests. Nothing in either run was mocked or replayed.
Chapter~5 reproduces the console evidence captured during both runs: the
pipeline mid-flight with live per-agent token counts, the fully green
timeline awaiting approval, the honest timed-out terminal state of the
first run, and the complete Approved-to-Succeeded timeline of the second.

\subsection{Run 1: Approval-Gate Timeout --- the Safety Mechanism Under Real Timing Pressure}

In the first run, a fault fixture was injected, Sentinel detected it, and
the full agent pipeline executed to completion: the Architect produced a
structured plan, the Backend agent generated a unified-diff patch, the QA
agent generated Hypothesis property tests, and the Validator executed the
patch against those tests inside its locked-down Docker sandbox --- every
Layer~1 and Layer~2 step reached a green state on the live timeline. The
candidate patch then arrived at the human approval gate, where it sat
pending. No human decision was made within the 120-second decision window,
the window elapsed, and the workflow resolved the incident to a rejected
terminal state.

It is important to read this outcome correctly: it is a \emph{positive}
test result, not a failure. The central safety invariant of the platform is
that no patch is ever deployed without passing validation and clearing the
approval gate, and that the system fails closed on ambiguity --- an absent
decision is treated as a refusal, never as consent. The first run exercised
exactly this invariant under real timing pressure, with a real generated
patch waiting on the other side of the gate, and the system behaved
precisely as designed: it discarded the work rather than shipping an
unreviewed change, and it recorded the timeout-rejected outcome honestly
and durably in the incident timeline. A demonstration environment could
easily have hidden this run and shown only the success; it is reported here
deliberately, because the rejected run is the stronger evidence that the
gate is real.

\subsection{Run 2: Approved and Succeeded in 1 Minute 40 Seconds}

The second run followed the same detection path but the pending approval
was acted on within the window. The incident was classified at Medium
severity, the generated patch passed sandbox validation, the approval was
granted through the console's decision dialog, and the pipeline proceeded
through its deploy stage to a \emph{Succeeded} terminal state. The complete
timeline --- from fault detection through every agent step to the final
succeeded state --- spanned 1~minute 40~seconds.

Two properties of this run make it strong evidence rather than a staged
demonstration. First, the accounting is real: the timeline records genuine
per-agent token consumption and cost figures for every LLM-backed step, as
produced by the live Ollama provider during the run. Second, the artifacts
are inspectable: the raw JSON payload of the Backend agent's activity event
shows the actual unified diff generated by \code{qwen2.5-coder:14b}, with
\code{provider="ollama"} recorded in the event metadata --- the patch that
was validated, approved, and deployed is the patch the model wrote, not a
canned fixture.

Together the two runs test both branches of the approval gate --- refusal
by timeout and explicit approval --- against the same real pipeline, which
is precisely the pair of behaviours a closed-loop recovery system must get
right.

\section{Live Deploy-Engine Verification}

The deploy-engine was verified with its own live end-to-end run against a
real public GitHub repository, \code{heroku/node-js-getting-started}, which
deliberately contains no \code{Dockerfile}. This exercises the engine's
hardest path: it must clone the repository, detect the stack from marker
files, generate a working \code{Dockerfile} from scratch, build the image,
run the container under resource limits, and prove liveness over HTTP.
Listing~\ref{lst:deploy-e2e} reproduces the actual request/response
transcript of the run.

\begin{lstlisting}[caption={Live deploy-engine end-to-end run: clone, generate Dockerfile, build, run, HTTP-verify, stop.},label={lst:deploy-e2e}]
POST /v1/deploy  (repo: heroku/node-js-getting-started, no Dockerfile present)
-> 200 OK
  "status": "running", "dockerfile_source": "generated", "detected_stack": "node",
  "url": "http://localhost:54671", "image_tag": "nexis-preview-test-project:22d06076c357"
  build_log: "...npm ci... EXPOSE 3000... CMD [\"npm\",\"start\"]..."
  container_log: "Listening on 3000\nRendering 'pages/index' for route '/'"

GET http://localhost:54671/  -> real rendered HTML of the Node app (verified)

POST /v1/deploy/test-e2e-001/stop -> {"status":"stopped"}; container removed.
\end{lstlisting}

The transcript demonstrates each stage genuinely occurring: the engine
detected a Node.js stack (\code{detected\_stack: "node"}), synthesised a
Dockerfile (\code{dockerfile\_source: "generated"}) whose build log shows
the generated \code{npm ci}, \code{EXPOSE 3000}, and \code{CMD} directives
executing; the container log shows the application itself booting and
serving its index route; an independent HTTP \code{GET} against the
returned preview URL returned the application's real rendered HTML; and the
stop endpoint cleanly terminated and removed the container. Because the
repository is public and the request sequence is three HTTP calls, this is
a reproducible integration test, not a one-off observation.

One environmental substitution should be noted for reproducibility: Docker
itself was unavailable in the evaluation environment, so
Podman~\cite{podman} was used as the container runtime through the same
Docker-compatible API~\cite{docker} that the deploy-engine and validator
target. The engine's code path is identical under both runtimes, and the
run above proves the Podman path works end-to-end. The clone in this run
used unauthenticated public-repository access; the authenticated GitHub App
installation-token path exists in the code but was not exercised in this
evaluation environment, as stated in the limitations.

\section{Defects Found and Fixed During This Engineering Cycle}
\label{sec:defects}

The verification process was not merely confirmatory: it surfaced two
genuine defects, both of which were root-caused and fixed during the cycle.
Table~\ref{tab:defects} summarises them.

\begin{table}[htbp]
\caption{Defects discovered by the verification process, with root causes, fixes, and verification.}
\label{tab:defects}
\centering
\footnotesize
\begin{tabular}{@{}p{3.0cm}p{4.3cm}p{3.3cm}p{4.0cm}@{}}
\toprule
\textbf{Defect} & \textbf{Root Cause} & \textbf{Fix} & \textbf{Verification} \\
\midrule
13 JWT-related tests silently failing & Token \code{exp} claims were validated against the wall clock instead of an injected test clock, so time-sensitive assertions depended on when the suite happened to run & One-line change: inject the test clock via \code{jwt.WithTimeFunc} & All 13 tests now pass deterministically; production behaviour verified unchanged, because the production code's default clock is already identical to the JWT library's own default \\
\addlinespace
Database migration failed to apply & \code{current\_role} is a reserved PostgreSQL keyword and was used unquoted as a column name in a migration & Column renamed to \code{existing\_role} & Migration applies cleanly; the Go test suites over the affected schema remain green \\
\bottomrule
\end{tabular}
\end{table}

Both defects illustrate why the verification loop earns its cost. The JWT
defect was \emph{silent}: the affected tests were failing without blocking
anything visible, which is exactly the failure mode that erodes a test
suite's value over time; it was found only because the full
\code{go test ./...} output was read rather than skimmed. The migration
defect was latent in a reserved-word collision that only manifests at
migration-apply time --- the kind of bug that static review of the SQL text
can easily miss.

For completeness, one further pre-existing bug was found during this cycle
but deliberately not fixed, as it is unrelated to the recovery pipeline and
out of scope: the workspace-onboarding progress poller requests
\code{/v1/workspaces/\{id\}/events} and receives a 404, although the
workspace itself provisions successfully server-side regardless. It is
recorded here, and in the limitations, rather than quietly omitted.

\section{Testing Summary}

The evidence base for NEXIS, stated plainly, is as follows. Static
verification is 100\% clean across every service in the platform: four Go
modules pass build, vet, and their full test suites with zero errors; both
Python services pass \code{pytest}, including the causal-inference
service's 38 tests; and the frontend passes type checking, linting, and a
full production build. On top of that foundation sit three live
integration runs observed directly against the real running system --- two
complete nine-agent recovery pipeline executions (one timeout-rejected,
demonstrating the fail-closed approval gate; one Approved and Succeeded in
1~minute 40~seconds with real per-agent token and cost accounting) and one
complete deploy-engine run against a real public repository (clone,
generated Dockerfile, build, run, HTTP-verified, stopped). None of these
runs were simulated, and the LLM behind the recovery runs was a real local
model, not a stub.

Equally important is what was \emph{not} tested. No formal load or
performance testing was performed: the platform has not been exercised
under concurrent multi-tenant traffic, no throughput or latency figures
under load were measured, and consequently none are claimed in this
report. The single measured end-to-end duration (1~minute 40~seconds) is
the observed time of one real run, not a statistical performance claim.
Load testing under realistic multi-tenant concurrency, along with a larger
sample of recovery runs from which meaningful MTTR distributions could be
drawn, is named explicitly as future work in Chapter~7.
